\documentclass[twocolumn]{aastex631}
\begin{document}
\title{The reionization ionizing-photon-budget shortfall for star-forming galaxies is not robust to systematics at z\~{}7 (optimistic systematic corner)}
\author{NebulaMind Lab (autonomous pipeline)}
\affiliation{NebulaMind Lab, \url{https://lab.nebulamind.net}}
\begin{abstract}
Reionization ionizing-photon-budget reconciliation at z\~{}7: star-forming galaxies require f\_esc=0.022 (+0.019/-0.010) to close the budget (Madau-Dickinson SFRD, log xi\_ion=25.5+/-0.15, clumping C=2-5, JWST-SFRD tail), versus indirect-proxy-inferred f\_esc=0.080 (+0.147/-0.051) from LzLCS O32/beta calibrations. Median delta(required-inferred)=-0.055 dex-frac (16-84\%: -0.201 to -0.002); 14\% of the systematic MC shows a shortfall. Conclusion: the budget CLOSES within the systematic, and the sign holds under both O32 and beta calibrations. The result is bounded by the xi\_ion x clumping x proxy-calibration systematic, not by statistics. Generated autonomously from public data (jwst) via the NebulaMind Lab runner: a bounded, reproducible, descriptive study.
\end{abstract}
\section{Introduction}
Introduction: Recent studies have highlighted a potential crisis in the reionization photon budget, suggesting that star-forming galaxies may not produce enough ionizing photons to account for the observed reionization [Muñoz2024]. This discrepancy has sparked interest in reconciling the photon budget using updated calibrations and models. Previous work by Park et al. [Park2022] demonstrated the importance of accurately modeling reionization processes, while Davies et al. [Davies2021] emphasized the need for increased ionizing sources to meet the demands of absorption-dominated reionization.
\section{Data and method}
Data and method: To address this issue, we adopt a literature-anchored budget calculation approach, utilizing the Madau \& Dickinson (2014) analytic fitting function for the cosmic star formation rate density (SFRD). We also incorporate published values for xi\_ion and O32/beta f\_esc proxy calibrations from LzLCS [Chisholm+22, Flury+22; Simmonds+24]. By systematically reconciling these literature values, we aim to determine the required escape fraction (f\_esc) for star-forming galaxies to close the reionization photon budget at z\~{}7.
\section{Result}
Result: Our analysis reveals that star-forming galaxies require an escape fraction of f\_esc=0.022 (+0.019/-0.010) to reconcile the reionization ionizing-photon-budget at z\~{}7, assuming a Madau-Dickinson SFRD, log xi\_ion=25.5+/-0.15, and clumping factor C=2-5. This value is lower than the indirect-proxy-inferred f\_esc=0.080 (+0.147/-0.051) derived from LzLCS O32/beta calibrations. The median difference between the required and inferred escape fractions is -0.055 dex-frac, with 14\% of systematic Monte Carlo simulations showing a shortfall.
\begin{figure}\centering\includegraphics[width=\columnwidth]{result.png}\caption{The reionization ionizing-photon-budget shortfall for star-forming galaxies is not robust to systematics at z\~{}7 (optimistic systematic corner)}\end{figure}
\section{Caveats}
Caveats: It is essential to acknowledge that our approach relies on automated selection and uncalibrated measurements from published literature, which may introduce biases and limitations. The accuracy of our result depends on the assumptions made in the adopted models and calibrations, such as the Madau-Dickinson SFRD and O32/beta f\_esc proxy calibrations. Additionally, our analysis does not account for potential systematic errors or uncertainties in the underlying data used to derive these literature values. Further investigation and refinement of these assumptions are necessary to strengthen the conclusions drawn from this study.
\section*{References}
\footnotesize
[Muoz2024] Muñoz et al. (2024). Reionization after JWST: a photon budget crisis?. bibcode 2024MNRAS.535L..37M \par [Park2022] Park et al. (2022). Calibrating excursion set reionization models to approximately conserve ionizing photons. bibcode 2022MNRAS.517..192P \par [Duncan2015] Duncan et al. (2015). Powering reionization: assessing the galaxy ionizing photon budget at z < 10. bibcode 2015MNRAS.451.2030D \par [Davies2021] Davies et al. (2021). The Predicament of Absorption-dominated Reionization: Increased Demands on Ionizing Sources. bibcode 2021ApJ...918L..35D \par [Madau2017] Madau (2017). Cosmic Reionization after Planck and before JWST: An Analytic Approach. bibcode 2017ApJ...851...50M

\end{document}
